\section{Introducción}

\subsection{Punto de partida de la Etapa 1}

La primera etapa dejó un resultado favorable, pero también un problema bastante
claro. Los tres modelos ligeros superaron la meta de macro-F1 de 0.85:
EfficientNet-B0 alcanzó 0.9146, ShuffleNetV2-x1.0 obtuvo 0.9030 y
EfficientNet-Lite0 llegó a 0.8951. La diferencia entre ellos no estaba tanto en
reconocer roya, hojas sanas o necrosis letal, sino en las clases con menos
ejemplos. \texttt{potassium\_deficiency} tenía solamente 266 imágenes y su F1
bajó hasta 0.62 con EfficientNet-B0, 0.54 con ShuffleNet y 0.49 con
EfficientNet-Lite0. El modelo podía acertar casi siempre en el conjunto
completo y, al mismo tiempo, fallar justo en la clase que más necesitaba
mejorar.

La matriz de confusión de aquella etapa mostró que las deficiencias de
nitrógeno, fósforo y potasio compartían parte de sus errores. También aparecía
una confusión menor entre mancha gris y tizón norteño. Las explicaciones LIME y
Grad-CAM aportaron otra advertencia: a veces la red miraba la hoja y aun así se
equivocaba, pero en imágenes de laboratorio también podía aprovechar un fondo
uniforme que no estaría presente en una parcela. La confianza alta tampoco
garantizaba una decisión correcta. Esos hallazgos hicieron insuficiente
limitarse a entrenar más épocas o escoger el mayor valor de accuracy.

Había además un problema de evaluación. El conjunto llamado \emph{test} había
sido consultado durante comparaciones, exportaciones y verificaciones
posteriores. Sus métricas siguen siendo útiles como registro histórico, pero ya
no representan una prueba ciega. Presentarlas otra vez como resultado final
habría dado una precisión metodológica que el experimento no tenía.

\subsection{Alcance de la Etapa 2}

La segunda etapa responde a esos puntos con un protocolo nuevo. Primero se
reconstruyó el corpus desde la publicación completa, se verificó cada archivo y
se retiraron duplicados exactos antes de dividir. Luego se congeló un holdout
nuevo de 5\,015 imágenes, sin emplearlo para elegir arquitectura,
hiperparámetros ni pesos del ensemble. La selección trabajó únicamente con
entrenamiento y validación; el holdout se abrió una vez al final.

La máquina disponible no expuso una GPU CUDA. Para no reemplazar un experimento
real por cifras estimadas, se congelaron backbones preentrenados en ImageNet y
se optimizó sobre sus embeddings una cabeza lineal de pérdida logística. Esto
no equivale al ajuste fino \emph{end-to-end} empleado en la Etapa 1. Sí permite
comparar cuatro representaciones modernas sobre las 33\,433 imágenes únicas,
ejecutar 30 trials de Optuna, probar soft voting, repetir cinco folds y evaluar
el holdout completo en CPU. La restricción cambia la pregunta experimental:
se mide qué tan separables son las nueve clases con cada representación
preentrenada y una cabeza común optimizada. El informe conserva esa distinción
cada vez que compara con el baseline histórico.

También se revisó el producto que rodea al clasificador. La aplicación móvil
captura o selecciona una foto, ejecuta un modelo TFLite local, calcula
confianza y rechazo fuera de dominio, guarda el diagnóstico sin depender de
conectividad y permite sincronizar correcciones. Su modelo embarcado pertenece
a la línea histórica EfficientNet-Lite0; no se sustituyó silenciosamente por el
modelo experimental de esta etapa. La demo incluida aquí ejecuta la selección
de Etapa 2 sobre una imagen real y deja una salida reproducible, mientras la
inspección de la aplicación confirma que existe una ruta móvil completa.

El objetivo no fue acumular arquitecturas. EfficientNet-B0 y ShuffleNet
continuaron por su papel en la etapa anterior; MobileNetV3 Small añadió una
alternativa diseñada para dispositivos; FastViT-T8 permitió contrastar una
arquitectura híbrida reciente. EfficientNet-B4, GhostNetV2 y otros nombres
registrados en el código no aparecen como resultados porque no tuvieron una
corrida completa en este protocolo. Separar lo implementado de lo evaluado
evita convertir el catálogo de modelos en evidencia que no existe.

La lectura final se concentra en cambios observables: cuánto aportó el tuning,
si los modelos cometieron errores distintos, si el ensemble justificó su costo,
qué variación hubo entre folds, qué ocurrió con potasio y qué brechas aparecieron
entre clases, fuentes, resolución, nitidez y entornos de captura. La última
parte traslada esos hallazgos a decisiones de uso. Doctor Maíz puede orientar
una inspección; no identifica con certeza una causa agronómica ni prescribe por
sí solo una aplicación de insumos.
